% !TEX program = xelatex
% MDPSA System Overview — standalone presentation
\documentclass[8pt,aspectratio=169]{beamer}

\usetheme[numbering=fraction,progressbar=frametitle]{metropolis}
\usepackage{fontspec}
\setmainfont{Times New Roman}
\usepackage{booktabs}
\usepackage{xcolor}
\usepackage{array}

\definecolor{mblue}{RGB}{0,82,147}
\definecolor{agentcol}{RGB}{0,100,60}

\title{MDPSA: Multi-Agent Dissolution--Precipitation Simulation Agent}
\subtitle{Automated thermodynamic modeling of selective metal recovery}
\author{Vahid Attari}
\date{2026}
\institute{NRCan CC-NRCan-078}

\begin{document}

\maketitle

%% ── Slide 1: Motivation ───────────────────────────────────────────────────────
\begin{frame}{Motivation}
  \begin{columns}[T]
    \begin{column}{0.48\textwidth}
      \begin{block}{Problem}
        \begin{itemize}
          \setlength{\itemsep}{5pt}
          \item Battery leachates (NMC532) contain Cu, Co, Ni as targets
                and Al, Fe, Mn as impurities.
          \item Selective precipitation via pH control requires a
                complexing modifier that retains target metals in solution
                while allowing impurities to precipitate.
          \item Identifying the best modifier log$\,K$ combination is
                a high-dimensional search problem.
        \end{itemize}
      \end{block}
    \end{column}
    \begin{column}{0.48\textwidth}
      \begin{block}{Solution}
        \begin{itemize}
          \setlength{\itemsep}{5pt}
          \item Automate the full cycle: literature search $\to$
                PHREEQC simulation $\to$ analysis $\to$ optimization.
          \item Use a LangGraph multi-agent pipeline so each stage
                is modular, reproducible, and auditable.
          \item Close the loop with Bayesian optimization to suggest
                improved log$\,K$ values without exhaustive enumeration.
        \end{itemize}
      \end{block}
    \end{column}
  \end{columns}
  \vspace{4pt}
  \begin{beamercolorbox}[rounded=true, shadow=false, wd=\textwidth]{block body}
    \tiny
    Experimental reference: NRCan CC-NRCan-078. Leachate composition: Cu, Co, Ni, Mn, Al, Fe in H$_2$SO$_4$ medium. Modifier dose 0.20 mol/kgw.
  \end{beamercolorbox}
\end{frame}

%% ── Slide 2: Pipeline Architecture ───────────────────────────────────────────
\begin{frame}{Pipeline Architecture}
  \begin{columns}[T]
    \begin{column}{0.55\textwidth}
      \small
      \renewcommand{\arraystretch}{1.35}
      \begin{tabular}{@{}lp{6.2cm}@{}}
        \toprule
        \textbf{Agent} & \textbf{Role} \\
        \midrule
        \textbf{ScoutAgent}     & Queries Claude LLM for NIST formation
                                  constants; writes \texttt{.phr} (PHREEQC
                                  reaction block) and \texttt{.yaml}
                                  (species tracking list). \\[4pt]
        \textbf{SimRunnerAgent} & Runs PHREEQC titration (0 $\to$ 1.0 equiv
                                  NaOH, 60 steps) for open, closed, and
                                  partial CO$_2$ scenarios; parses output
                                  CSVs. \\[4pt]
        \textbf{AnalystAgent}   & Computes dissolved-metal fractions, RMSE
                                  vs. experiment, selectivity scorecard, and
                                  species-contribution figures. \\[4pt]
        \textbf{BOAgent}        & Gaussian-process Bayesian optimization
                                  over log$\,K$ space; proposes next
                                  modifier candidate. \\[4pt]
        \textbf{BeamerAgent}    & Assembles and compiles this XeLaTeX PDF
                                  automatically at end of each run. \\
        \bottomrule
      \end{tabular}
    \end{column}
    \begin{column}{0.42\textwidth}
      \small
      \begin{block}{Orchestration}
        \begin{itemize}
          \setlength{\itemsep}{4pt}
          \item \textbf{Framework}: LangGraph \texttt{StateGraph}
          \item \textbf{State}: typed \texttt{SimState} dict passed
                between nodes
          \item \textbf{Routing}: conditional edges; SimRunner retries
                on PHREEQC error (max 3)
          \item \textbf{LLM}: Anthropic Claude (claude-sonnet-4-5)
                via Messages API
        \end{itemize}
      \end{block}
      \vspace{4pt}
      \begin{block}{Flow}
        \centering
        \small
        Scout $\to$ SimRunner $\to$ Analyst $\to$ BO $\to$ Beamer
      \end{block}
    \end{column}
  \end{columns}
  \vspace{4pt}
  \begin{beamercolorbox}[rounded=true, shadow=false, wd=\textwidth]{block body}
    \tiny
    Modifier files in \texttt{database/modifiers/}; simulation outputs in \texttt{experiments/single/}; figures and PDFs in \texttt{presentations/\{modifier\}/}.
  \end{beamercolorbox}
\end{frame}

%% ── Slide 3: Geochemical Model ────────────────────────────────────────────────
\begin{frame}{Geochemical Model}
  \begin{columns}[T]
    \begin{column}{0.48\textwidth}
      \begin{block}{PHREEQC setup}
        \begin{itemize}
          \setlength{\itemsep}{4pt}
          \item Thermodynamic database: \texttt{sit.dat} (SIT activity model)
          \item NaOH titration: 0 $\to$ 1.0 equiv in 60 kinetic steps
          \item Equilibrium phases allowed to precipitate:
                Gibbsite, Ferrihydrite(am), Cu(OH)$_2$(s),
                Co(OH)$_2$(s), Ni(OH)$_2$(s), MnOOH
          \item Three CO$_2$ scenarios: open (atmosphere), closed
                (sealed), partial (0.1 atm)
          \item Fe complexation suppressed (kinetic limitation)
        \end{itemize}
      \end{block}
    \end{column}
    \begin{column}{0.48\textwidth}
      \begin{block}{Complex types in \texttt{.phr}}
        \small
        \renewcommand{\arraystretch}{1.3}
        \begin{tabular}{@{}lll@{}}
          \toprule
          \textbf{Type} & \textbf{Reaction} & \textbf{Example} \\
          \midrule
          ML    & M$^{n+}$ + L = ML    & CuEdta$^{2-}$ \\
          ML$_2$ & M$^{n+}$ + 2L = ML$_2$ & CuEdta$_2^{6-}$ \\
          MOHL  & M + L $-$ H$^+$ + H$_2$O = product & CuOHEdta \\
          \bottomrule
        \end{tabular}
        \vspace{6pt}
        \begin{itemize}
          \setlength{\itemsep}{3pt}
          \item log$\,K$ values from NIST SRD 46 v8.0
          \item Values tuned iteratively against experiment
          \item BO closes the tuning loop automatically
        \end{itemize}
      \end{block}
    \end{column}
  \end{columns}
  \vspace{4pt}
  \begin{beamercolorbox}[rounded=true, shadow=false, wd=\textwidth]{block body}
    \tiny
    NIST Standard Reference Database 46 Version 8.0. SIT: Specific Ion Interaction Theory. Temperature: 25$^\circ$C.
  \end{beamercolorbox}
\end{frame}

%% ── Slide 4: Selectivity Scoring ──────────────────────────────────────────────
\begin{frame}{Selectivity Scoring and Optimization}
  \begin{columns}[T]
    \begin{column}{0.48\textwidth}
      \begin{block}{Scorecard logic}
        \begin{itemize}
          \setlength{\itemsep}{4pt}
          \item Comparison points: 0.3 and 1.0 equiv NaOH
          \item \textbf{Target metals} (Cu, Co, Ni): pass if dissolved
                fraction $> 70\%$ at each comparison point
          \item \textbf{Impurity metals} (Al, Fe, Mn): pass if dissolved
                fraction $< 30\%$
          \item Score: 0, 1, or 2 per metal (one point per equiv point)
          \item RMSE computed over all experimental equiv points
        \end{itemize}
      \end{block}
    \end{column}
    \begin{column}{0.48\textwidth}
      \begin{block}{Bayesian optimization}
        \begin{itemize}
          \setlength{\itemsep}{4pt}
          \item Objective: maximize total selectivity score
          \item Search space: log$\,K$ of each complex $\pm 3$ units
                around NIST reference
          \item Surrogate: Gaussian process with Mat\'{e}rn kernel
          \item Acquisition: expected improvement (EI)
          \item Each pipeline run adds one observation; BO proposes
                the next log$\,K$ set
        \end{itemize}
      \end{block}
    \end{column}
  \end{columns}
  \vspace{4pt}
  \begin{beamercolorbox}[rounded=true, shadow=false, wd=\textwidth]{block body}
    \tiny
    BO implemented with \texttt{scikit-learn} GaussianProcessRegressor. Modifier dose and leachate composition fixed; only log$\,K$ values are optimized.
  \end{beamercolorbox}
\end{frame}

%% ── Slide 5: Outputs and Reproducibility ──────────────────────────────────────
\begin{frame}{Outputs and Reproducibility}
  \begin{columns}[T]
    \begin{column}{0.48\textwidth}
      \begin{block}{Per-run outputs}
        \small
        \renewcommand{\arraystretch}{1.3}
        \begin{tabular}{@{}ll@{}}
          \toprule
          \textbf{File} & \textbf{Content} \\
          \midrule
          \texttt{.phr}           & PHREEQC reaction block \\
          \texttt{.yaml}          & Species tracking list \\
          \texttt{neutralization.csv} & Titration results \\
          \texttt{*\_model\_vs\_exp.png} & Model vs. experiment \\
          \texttt{*\_species\_contribution.png} & Species breakdown \\
          \texttt{*\_summary.pdf} & This presentation \\
          \texttt{run\_config.yaml} & Full run history \\
          \bottomrule
        \end{tabular}
      \end{block}
    \end{column}
    \begin{column}{0.48\textwidth}
      \begin{block}{Reproducibility}
        \begin{itemize}
          \setlength{\itemsep}{4pt}
          \item All modifier files version-controlled in Git
          \item \texttt{run\_config.yaml} logs every run with
                scorecard, figure paths, and timestamp
          \item \texttt{--force-scout} flag regenerates LLM output;
                omitting it reuses existing \texttt{.phr}/\texttt{.yaml}
          \item Presentations committed to
                \texttt{presentations/\{modifier\}/}
          \item \texttt{.tex} files Overleaf-ready
                (\texttt{\% !TEX program = xelatex})
        \end{itemize}
      \end{block}
    \end{column}
  \end{columns}
  \vspace{4pt}
  \begin{beamercolorbox}[rounded=true, shadow=false, wd=\textwidth]{block body}
    \tiny
    Repository: \texttt{github.com/vahid2364/modifier-driven-precipitation-simulation-agent}. Stack: Python 3.11, LangGraph, PHREEQC v3, scikit-learn, XeLaTeX (TeX Live 2026).
  \end{beamercolorbox}
\end{frame}

\end{document}
